% Content shared by the web and print versions.
%
% Each statement carries:
%   \lean{Name, Name, ...}   the Lean declaration(s) that formalise it
%   \leanok                  the statement (or, inside a proof, the proof) is formalised
%   \uses{label, ...}        dependencies, which become edges of the dependency graph
% Statements without \leanok show up in the dependency graph as not yet
% formalised. Two such placeholders are included below on purpose.
%
% A proof must come AFTER \end{theorem} (not inside it): the dependency-graph
% plugin attaches a proof to the environment that precedes it. Only then do
% the proof's \uses become graph edges and its \leanok fill the node green.

\chapter{Introduction}

This blueprint links the statements of \emph{Partitioning 3-homogeneous
latin bitrades} (arXiv:0710.0938) to their formal counterparts in the Lean 4
project \texttt{Partitioning3Homogeneous}. Every numbered statement below
names the Lean declaration that formalises it, and the dependency graph
records which statements are used to prove which.

Two conventions differ between the paper and Lean.
\begin{itemize}
  \item The paper lets permutations act on the right, so $x\tau_1\tau_2$
    means ``apply $\tau_1$, then $\tau_2$''. Lean's \texttt{Equiv.Perm}
    composes right to left, so the same composite is written
    \texttt{tau2 * tau1}. In particular (T1), $\tau_1\tau_2\tau_3 = 1$, is
    the Lean statement \texttt{tauSymbol * tauColumn * tauRow = 1}.
  \item The paper works with $n \times n$ arrays. Lean works with arbitrary
    finite index types for rows, columns, and symbols, and 3-homogeneity is
    required of every element of those types. A bitrade with empty rows must
    be re-indexed before the main theorem applies.
\end{itemize}

\chapter{Latin bitrades}

\begin{definition}[Partial Latin square]\label{def:partial-latin}
  \lean{LatinBitrade.Entry, LatinBitrade.IsPartialLatin}\leanok
  A \emph{partial Latin square} is a set $T \subseteq A_1 \times A_2 \times A_3$
  of entries $(i, j, k)$, read as ``symbol $k$ in row $i$ and column $j$'',
  such that two entries agreeing in any two coordinates agree in the third.
\end{definition}

\begin{definition}[Latin bitrade; Definition 2.1]\label{def:bitrade}
  \lean{LatinBitrade.HasUniqueMates, LatinBitrade.IsBitrade}\leanok
  \uses{def:partial-latin}
  Let $\Tpos, \Tneg$ be partial Latin squares. The pair $(\Tpos, \Tneg)$
  is a \emph{bitrade} if
  \begin{itemize}
    \item[(R1)] $\Tpos \cap \Tneg = \emptyset$;
    \item[(R2)] for all $(a_1, a_2, a_3) \in \Tpos$ and all $r \neq s$ there is
      a unique $(b_1, b_2, b_3) \in \Tneg$ with $a_r = b_r$ and $a_s = b_s$;
    \item[(R3)] the same with $\Tpos$ and $\Tneg$ exchanged.
  \end{itemize}
  In Lean, (R2) and (R3) are \texttt{HasUniqueMates positive negative} and
  \texttt{HasUniqueMates negative positive}.
\end{definition}

\begin{definition}[$k$-homogeneous]\label{def:homogeneous}
  \lean{LatinBitrade.IsKHomogeneous}\leanok
  A bitrade is \emph{$k$-homogeneous} if every row and every column of
  $\Tpos$ contains $k$ entries and every symbol occurs $k$ times in $\Tpos$.
\end{definition}

\begin{definition}[Transversal]\label{def:transversal}
  \lean{LatinBitrade.IsTransversal, LatinBitrade.IsThreeTransversalPartition}\leanok
  A set $X \subseteq \Tpos$ is a \emph{transversal} if it meets each row and
  each column of $\Tpos$ in exactly one entry and
  $\lvert \{ k \mid (i, j, k) \in X \} \rvert = \lvert X \rvert$.
  \texttt{IsThreeTransversalPartition} packages three transversals which
  are pairwise disjoint and cover $\Tpos$.
\end{definition}

\begin{definition}[The mate bijections $\beta_r$]\label{def:beta}
  \lean{LatinBitrade.Canonical.beta₁, LatinBitrade.Canonical.beta₂, LatinBitrade.Canonical.beta₃}\leanok
  \uses{def:bitrade}
  For $r \in \{1, 2, 3\}$ define $\beta_r \colon \Tneg \to \Tpos$ by
  $(a_1, a_2, a_3)\beta_r = (b_1, b_2, b_3)$ if and only if $a_r \neq b_r$
  and $a_i = b_i$ for $i \neq r$. By (R2) and (R3) each $\beta_r$ is a
  bijection. In Lean the bijections are built from the unique-mate clauses
  with \texttt{mateEquivBy}; the inequality $a_r \neq b_r$ follows from (R1).
\end{definition}

\begin{definition}[The $\tau_i$ representation; Equation (1)]\label{def:tau}
  \lean{LatinBitrade.Canonical.tauRow, LatinBitrade.Canonical.tauColumn, LatinBitrade.Canonical.tauSymbol}\leanok
  \uses{def:beta}
  Define permutations of $\Tpos$ by
  \[
    \tau_1 = \beta_2^{-1}\beta_3, \qquad
    \tau_2 = \beta_3^{-1}\beta_1, \qquad
    \tau_3 = \beta_1^{-1}\beta_2,
  \]
  maps applied from left to right. The cycles of $\tau_1$, $\tau_2$, $\tau_3$
  are the rows, columns, and symbols of $\Tpos$.
\end{definition}

\begin{lemma}[(T1) and coordinate preservation]\label{lem:T1}
  \lean{LatinBitrade.Canonical.tau_product, LatinBitrade.Canonical.tauRow_preserves_row, LatinBitrade.Canonical.tauColumn_preserves_column, LatinBitrade.Canonical.tauSymbol_preserves_symbol}\leanok
  \uses{def:tau}
  $\tau_1\tau_2\tau_3 = 1$, and $\tau_1$ preserves rows, $\tau_2$ preserves
  columns, $\tau_3$ preserves symbols.
\end{lemma}
\begin{proof}\leanok
  Unfold the definitions; consecutive $\beta$'s cancel. Each $\tau_i$ is a
  composite of two maps that both fix the $i$th coordinate.
\end{proof}

\begin{lemma}[(T3): fixed-point-freeness]\label{lem:T3}
  \lean{LatinBitrade.Canonical.tauRow_fixedPointFree, LatinBitrade.Canonical.tauColumn_fixedPointFree, LatinBitrade.Canonical.tauSymbol_fixedPointFree}\leanok
  \uses{def:tau, def:bitrade}
  Each $\tau_i$ is fixed-point-free.
\end{lemma}
\begin{proof}\leanok
  If $x\tau_1 = x$ then $y = x\beta_2^{-1} \in \Tneg$ satisfies
  $y\beta_2 = y\beta_3 = x$, so $y$ agrees with $x$ in all three
  coordinates and $x \in \Tpos \cap \Tneg$, contradicting (R1).
\end{proof}

\begin{lemma}[Order three; Equation (4)]\label{lem:order-three}
  \lean{LatinBitrade.perm_pow_three_of_fiber_card_three, LatinBitrade.Canonical.tauRow_pow_three, LatinBitrade.Canonical.tauColumn_pow_three, LatinBitrade.Canonical.tauSymbol_pow_three}\leanok
  \uses{lem:T1, lem:T3, def:homogeneous}
  If $(\Tpos, \Tneg)$ is 3-homogeneous then $\tau_i^3 = 1$ for
  $i = 1, 2, 3$.
\end{lemma}
\begin{proof}\leanok
  $\tau_1$ permutes the three entries of each row without fixed points,
  so it acts on each row as a 3-cycle. The Lean proof first rules out a
  2-cycle inside a 3-element fibre, then shows $x, x\tau_1, x\tau_1^2$
  exhaust the fibre.
\end{proof}

\begin{theorem}[Drápal; Theorem 2.4]\label{thm:drapal}
  A bitrade is equivalent up to isotopism to three permutations satisfying
  (T1), (T2), (T3), and it is primary if (T4) also holds.
  \emph{Not formalised: the main theorem does not need it. Condition (T2)
  is checked only for the worked example.}
\end{theorem}

\chapter{The Euclidean triangle group}

\begin{definition}[The group $\Gamma$; Equation (3)]\label{def:gamma}
  \lean{EuclideanTriangleGroup.TriangleGroup, EuclideanTriangleGroup.rho₁, EuclideanTriangleGroup.rho₂, EuclideanTriangleGroup.rho₃}\leanok
  The paper takes
  $\Gamma = \langle \rho_1, \rho_2, \rho_3 \mid
    \rho_1^3 = \rho_2^3 = \rho_3^3 = \rho_1\rho_2\rho_3 = 1 \rangle$,
  the orientation-preserving $(3,3,3)$ triangle group acting on the
  triangular tessellation of the plane.
  Lean uses the concrete model $\Z[\omega] \rtimes \Z/3$, with
  $\Z[\omega] \cong \Z^2$ in the basis $(1, \omega)$ and $\Z/3$ acting by
  multiplication by $\omega$. The generators are the affine maps
  \[
    \rho_1 \colon z \mapsto \omega z, \qquad
    \rho_2 \colon z \mapsto \omega z + 1, \qquad
    \rho_3 \colon z \mapsto \omega z + 1 + \omega,
  \]
  anticlockwise rotations by $2\pi/3$ about the three vertices
  $0$, $(2 + \omega)/3$, $(1 + 2\omega)/3$ of one triangle.
\end{definition}

\begin{lemma}[Relations]\label{lem:rho-relations}
  \lean{EuclideanTriangleGroup.rho_relations}\leanok
  \uses{def:gamma}
  $\rho_1^3 = \rho_2^3 = \rho_3^3 = \rho_1\rho_2\rho_3 = 1$ in the concrete
  model.
\end{lemma}
\begin{proof}\leanok
  By computation (\texttt{decide}).
\end{proof}

\begin{lemma}[Rotations are conjugate to generators]\label{lem:conjugate}
  \lean{EuclideanTriangleGroup.conjugate_rotation_one, EuclideanTriangleGroup.rotation_one_conjugate_generator}\leanok
  \uses{def:gamma}
  Every element of $\Gamma$ whose rotational component is $\omega$ is
  conjugate by a translation to one of $\rho_1$, $\rho_2$, $\rho_3$.
  This is the algebraic replacement for Cases 1--4 in the proof of
  Lemma 4.3 of the paper.
\end{lemma}
\begin{proof}\leanok
  Conjugating $z \mapsto \omega z + (m + n\omega)$ by the translation
  by $p + q\omega$ gives $z \mapsto \omega z + (m + p + q) + (n + 2q - p)\omega$.
  The residue $m + n \bmod 3$ is invariant and equals $0, 1, 2$ on
  $\rho_1, \rho_2, \rho_3$; choose $p, q$ to reach the representative.
\end{proof}

\begin{proposition}[The representation $\theta$; Equation (4)]\label{prop:theta}
  \lean{EuclideanTriangleGroup.generator_relations_force_translations_commute, EuclideanTriangleGroup.triangleRepresentation, EuclideanTriangleGroup.triangleRepresentation_rho₁, EuclideanTriangleGroup.triangleRepresentation_rho₂, EuclideanTriangleGroup.triangleRepresentation_rho₃}\leanok
  \uses{def:gamma}
  Let $a, b, c$ be elements of a group $G$ with
  $a^3 = b^3 = c^3 = abc = 1$. There is a homomorphism
  $\theta \colon \Gamma \to G$ with $\rho_1 \mapsto a$, $\rho_2 \mapsto b$,
  $\rho_3 \mapsto c$.
\end{proposition}
\begin{proof}\leanok
  Put $u = ba^{-1}$ and $v = aua^{-1}$. From $b^3 = 1$ one gets
  $ava^{-1} = (uv)^{-1}$, and from $c^3 = 1$ that $u$ and $v$ commute.
  So $(u, v)$ is a lattice representation on which conjugation by $a$
  acts as multiplication by $\omega$, and the universal property of the
  semidirect product gives $\theta$.
\end{proof}

\chapter{The three-colouring}

\begin{lemma}[Equivariant orbit maps; Lemma 4.1]\label{lem:orbit-map}
  \lean{EquivariantOrbitMap.MapsStabilizer, EquivariantOrbitMap.exists_equivariant}\leanok
  Let $\theta \colon \Gamma \to G$ be a homomorphism, $X$ a transitive
  $\Gamma$-set with base point $x_0$, and $Y$ a $G$-set with base point
  $y_0$. If $\theta$ maps the stabiliser of $x_0$ into the stabiliser of
  $y_0$, then $\gamma \cdot x_0 \mapsto \theta(\gamma) \cdot y_0$ is a
  well-defined equivariant map $X \to Y$.
  The stabiliser hypothesis is exactly what the paper's proof of Lemma 4.1
  leaves implicit; there it holds because $\Gamma$ acts freely on the
  shaded triangles.
\end{lemma}
\begin{proof}\leanok
  If $\gamma_1 x_0 = \gamma_2 x_0$ then $\gamma_2^{-1}\gamma_1$ fixes
  $x_0$, so $\theta(\gamma_2^{-1}\gamma_1)$ fixes $y_0$ and the two
  candidate values agree. Equivariance is the same computation.
\end{proof}

\begin{lemma}[Stabilisers are translations; Lemma 4.3]\label{lem:stabilizer}
  \lean{EuclideanTriangleGroup.rotation_one_not_in_stabilizer, EuclideanTriangleGroup.stabilizer_has_trivial_rotation}\leanok
  \uses{lem:conjugate}
  Let $\Gamma$ act on a set $D$ so that $\rho_1$, $\rho_2$, $\rho_3$ have
  no fixed points. Then every point stabiliser is contained in the
  translation subgroup $\Z[\omega]$.
\end{lemma}
\begin{proof}\leanok
  An element with rotational component $\omega$ fixing $x$ is conjugate
  by a translation $t$ to some $\rho_i$, which then fixes $t \cdot x$.
  The component $\omega^2$ case follows by taking inverses.
\end{proof}

\begin{theorem}[Three-colouring; Definition 4.2 and Corollary 4.5]\label{thm:coloring}
  \lean{EuclideanTriangleGroup.exists_three_coloring_of_transitive, EuclideanTriangleGroup.exists_three_coloring, EuclideanTriangleGroup.permutations_admit_three_coloring}\leanok
  \uses{lem:orbit-map, lem:stabilizer, prop:theta, lem:rho-relations}
  Let $a, b, c$ be fixed-point-free permutations of a set $D$ with
  $a^3 = b^3 = c^3 = abc = 1$. Then there is a colouring
  $\chi \colon D \to \Z/3$ with
  $\chi(ax) = \chi(bx) = \chi(cx) = \chi(x) + 1$ for all $x$.
  Neither (T2) nor transitivity (T4) is needed: the colouring is built on
  each orbit separately.
\end{theorem}
\begin{proof}\leanok
  Let $\Gamma$ act on $D$ through $\theta$. On the orbit of $x_0$ define
  $\chi(\gamma \cdot x_0)$ to be the rotational component of $\gamma$.
  By Lemma~\ref{lem:stabilizer} the stabiliser of $x_0$ has trivial
  rotational component, so Lemma~\ref{lem:orbit-map} applied to the
  projection $\Gamma \to \Z/3$ shows $\chi$ is well defined. Each $\rho_i$
  has rotational component $1$, so it advances $\chi$ by one.
\end{proof}

\chapter{Transversals}

\begin{lemma}[Colour classes are transversals; Lemma 4.6]\label{lem:classes}
  \lean{LatinBitrade.CyclesFibers, LatinBitrade.cyclesFibers_of_card_three, LatinBitrade.color_class_fiber_card, LatinBitrade.three_homogeneous_permutation_partition}\leanok
  \uses{thm:coloring, def:transversal}
  Let $\tau_1, \tau_2, \tau_3$ be fixed-point-free permutations of order
  three of a finite set $D$ with $\tau_3\tau_2\tau_1 = 1$ (Lean order),
  preserving coordinates $\mathrm{row}$, $\mathrm{col}$, $\mathrm{sym}$
  whose fibres all have three elements. Then the three colour classes of
  Theorem~\ref{thm:coloring} each meet every fibre of every coordinate
  exactly once, are pairwise disjoint, and cover $D$.
\end{lemma}
\begin{proof}\leanok
  A fibre of $\mathrm{row}$ is $\{x, x\tau_1, x\tau_1^2\}$, whose colours
  are $\chi(x), \chi(x) + 1, \chi(x) + 2$. So each colour occurs exactly
  once per fibre. The same holds for columns and symbols.
\end{proof}

\begin{lemma}[Transfer to entries]\label{lem:transfer}
  \lean{LatinBitrade.forgetDarts_isTransversal, LatinBitrade.coordinatePartition_isThreeTransversalPartition, LatinBitrade.exists_three_transversal_partition_from_permutations}\leanok
  \uses{lem:classes, def:transversal, def:homogeneous}
  A partition of the entries of $\Tpos$, regarded as an abstract dart set,
  into three coordinate transversals is a partition of $\Tpos$ into three
  transversals in the sense of Definition~\ref{def:transversal}.
\end{lemma}
\begin{proof}\leanok
  Forget the membership proofs. The symbol-count clause follows from
  injectivity of the symbol map on each class.
\end{proof}

\begin{theorem}[Cavenagh; Theorem 1.1]\label{thm:main}
  \lean{LatinBitrade.theorem_1_1}\leanok
  \uses{def:bitrade, def:homogeneous, def:transversal}
  Let $(\Tpos, \Tneg)$ be a 3-homogeneous bitrade. Then $\Tpos$ can be
  partitioned into three transversals.
\end{theorem}
\begin{proof}\leanok
  \uses{lem:T1, lem:T3, lem:order-three, lem:transfer}
  Apply Lemma~\ref{lem:transfer} to $\tau_1, \tau_2, \tau_3$ of
  Definition~\ref{def:tau}. The hypotheses are supplied by
  Lemmas~\ref{lem:T1}, \ref{lem:T3}, and \ref{lem:order-three}.
\end{proof}

\begin{lemma}[Tessellation on the torus; Lemma 3.5]\label{lem:torus}
  \uses{def:tau}
  A 3-homogeneous bitrade defines a tessellation of the Euclidean plane by
  shaded and unshaded triangles, whose quotient is a torus.
  \emph{Not formalised: the geometric picture is replaced by
  Definition~\ref{def:gamma} and Lemma~\ref{lem:stabilizer}.}
\end{lemma}

\chapter{The worked example}

\begin{proposition}[Example 2.6 and its transversals]\label{prop:example}
  \lean{LatinBitrade.PaperExample.equation2_is_bitrade, LatinBitrade.PaperExample.equation2_is_three_homogeneous, LatinBitrade.PaperExample.displayed_partition_is_valid}\leanok
  \uses{def:bitrade, def:homogeneous, def:transversal}
  The two arrays of Equation (2) form a 3-homogeneous bitrade, and the
  three sets displayed after Lemma 3.5 partition $\Tpos$ into
  transversals.
\end{proposition}
\begin{proof}\leanok
  Kernel computation on the twelve entries (\texttt{decide}).
\end{proof}

\begin{proposition}[The printed $\tau_i$ satisfy (T1)--(T4)]\label{prop:example-tau}
  \lean{LatinBitrade.PaperExample.tau_equation_from_betas, LatinBitrade.PaperExample.tau_product_identity, LatinBitrade.PaperExample.tau_cycles_pairwise_disjoint, LatinBitrade.PaperExample.tau_fixedPointFree, LatinBitrade.PaperExample.tau_action_transitive}\leanok
  \uses{prop:example, def:tau}
  The permutations printed after Equation (2) are those of Equation (1),
  and they satisfy (T1), (T2), (T3), and (T4).
\end{proposition}
\begin{proof}\leanok
  Kernel computation. Transitivity is certified by an explicit word in
  the generators for every ordered pair of darts.
\end{proof}
